\section{Electrical control in fixed regions}\label{sec:network}

We next obtain electrical estimates in the reference coordinate.  The
estimates concern a fixed annulus or a fixed route between open sets.  Their
constants depend on relative geometry, while the discrete index from which
they hold may depend on the region and on the sample.  This is the
macroscopic control used to identify the limiting energy.

The canonical companion proves nondegeneration of extremal length on its
piecewise Euclidean flag surface.  We require the corresponding statement
for paths on the original graph edges.  We prove it here using the aligned
annuli of \cite[Theorem~6.4]{C} and the numerical observations of
Section~\ref{sec:locality}.  The density and path compactness arguments
follow the mechanism in \cite[Section~8]{C}.  For edge paths, the final
obstruction is the intersection of paired primal and dual edges.

\subsection{Edge extremal length and traffic}

Let \(\mathsf G\) be a finite undirected multigraph.  Its edges are always
counted as distinct occurrences.  For a finite walk \(\pi\), let
\(N_e(\pi)\) be the number of traversals of \(e\), in either direction and
with multiplicity.  For a nonnegative edge function \(g\), set
\[
 L_g(\pi)=\sum_{e\in E(\mathsf G)}g_eN_e(\pi),
 \qquad
 \norm{g}_2^2=\sum_{e\in E(\mathsf G)}g_e^2.
\]
Choose one orientation of each edge occurrence, write
\(d_eu=u(e^+)-u(e^-)\), and set
\[
 \mathcal E_{\mathsf G}^0(u)
   =\sum_{e\in E(\mathsf G)}(d_eu)^2.
\]
The choice of orientation is immaterial, and a loop has zero difference.
For a nonempty family \(\mathcal F\) of such walks, define
\begin{equation}\label{eq:edge-el}
 \EL_{\mathsf G}(\mathcal F)
 =\sup_{\substack{g\in[0,\infty)^{E(\mathsf G)}\\g\ne0}}
       \frac{\bigl(\inf_{\pi\in\mathcal F}L_g(\pi)\bigr)^2}
            {\norm{g}_2^2}.
\end{equation}
We set \(\EL_{\mathsf G}(\varnothing)=\infty\).  If the graph has no edges
and the family contains the empty walk, its extremal length is \(0\).
For a probability law \(\mu\) on walks, its traffic and squared traffic
are
\begin{equation}\label{eq:edge-traffic}
 \tau_\mu(e)=\EE_\mu N_e(\pi),
 \qquad
 \mathcal T_{\mathsf G}(\mu)=\sum_{e\in E(\mathsf G)}\tau_\mu(e)^2.
\end{equation}
These definitions, like \(\mathcal E^0\), have no factor \(1/2\).

For a finite family of objects with nonzero prescribed edge usages, the
probabilistic variational formula below is the reciprocal form of
\cite[Theorem~5.1]{AlbinPoggiCorradini2016}, since our \(\EL\) equals
their \(\operatorname{Mod}_2^{-1}\).  We include a proof using a
closed convex hull that also permits an arbitrary family of finite walks.  The
electrical convention and the port identity are the unit-conductance
case of \cite[Propositions~2.2--2.5]{BinderPechersky2025}, where the
extremal-length identity is attributed to Duffin.

\begin{lemma}[Finite variational identities]\label{lem:finite-el}
For a nonempty family of finite walks,
\begin{equation}\label{eq:el-traffic-identity}
 \EL_{\mathsf G}(\mathcal F)
 =\inf_{\mu}\mathcal T_{\mathsf G}(\mu),
\end{equation}
where finite-support probability laws on \(\mathcal F\) suffice for the
infimum.  General probability laws with finite expected edge counts give the
same infimum.

Let \(A,B\) be nonempty disjoint vertex sets, and let
\(\mathcal F_{\mathsf G}(A,B)\) contain all walks in \(\mathsf G\) from
\(A\) to \(B\).  Then
\begin{equation}\label{eq:port-resistance}
 \EL_{\mathsf G}(\mathcal F_{\mathsf G}(A,B))
 =R_{\mathrm{eff}}^{\mathsf G}(A,B)
 =\left[
   \inf_{\substack{u|_A=0\\u|_B=1}}\mathcal E_{\mathsf G}^0(u)
   \right]^{-1}.
\end{equation}
The two port sets are wired separately, and the resistance is computed in
the same graph \(\mathsf G\) that confines the walks.  If there is no
connecting walk, the energy infimum is \(0\) and its reciprocal is
\(\infty\).
\end{lemma}

\begin{proof}
Write \(a_\pi=(N_e(\pi))_{e\in E(\mathsf G)}\), and let \(K\) be the closed
convex hull of these nonnegative vectors.  For \(t\in K\) and \(g\ge0\),
\[
 \inf_{\pi\in\mathcal F}a_\pi\cdot g
 \le t\cdot g\le\norm{t}_2\norm{g}_2.
\]
Consequently \(\EL_{\mathsf G}(\mathcal F)\le\inf_{t\in K}\norm{t}_2^2\).
Let \(p\) be the point of \(K\) nearest \(0\).  If \(p=0\), this already
gives equality.  Otherwise \(p\ge0\), and the projection inequality gives
\[
 a_\pi\cdot p\ge\norm{p}_2^2
 \quad\hbox{for every }\pi\in\mathcal F.
\]
Taking the admissible nonzero density \(g=p\) in \eqref{eq:edge-el} gives
the reverse inequality.  Finite convex combinations are precisely the
traffic vectors of finite-support laws and are dense in \(K\).  A general
traffic vector with finite coordinates belongs to \(K\), by approximating
its law by finite laws and allowing an arbitrarily small error in the
coordinates.  This proves \eqref{eq:el-traffic-identity}.

For the port identity, first suppose a connecting walk exists.  Let
\(u|_A=0\), \(u|_B=1\), and take
\(g_e=\abs{d_eu}\).  A connecting walk and these port values force
\(g\ne0\), so this density is admissible in \eqref{eq:edge-el}.
The telescoping inequality gives \(L_g(\pi)\ge1\) on
every connecting walk, so \(\EL\ge1/\mathcal E_{\mathsf G}^0(u)\).
Conversely, for a density with
\(b=\inf_{\pi\in\mathcal F_{\mathsf G}(A,B)}L_g(\pi)>0\), let \(d_g(A,v)\)
be weighted graph distance from \(A\), with value \(\infty\) in a component
not meeting \(A\).  The function
\[
 u(v)=\min\{1,d_g(A,v)/b\}
\]
has the required port values and satisfies
\(\abs{d_eu}\le g_e/b\) on every edge.  Thus the energy infimum in
\eqref{eq:port-resistance} is at most \(\norm g_2^2/b^2\).  Take the
supremum over \(g\).  If no connecting walk exists, assign a constant to
each component, proving the stated convention.  The equality with wired
effective resistance is the Dirichlet principle for a unit-conductance
network.  Loop differences vanish throughout.
\end{proof}

We will repeatedly use the following consequence of the definition.  If
every walk of a family \(\mathcal F\) contains a subwalk in
\(\mathcal G\), then
\begin{equation}\label{eq:el-subpassage}
 \EL(\mathcal F)\ge\EL(\mathcal G),
\end{equation}
when the densities for \(\mathcal G\) are extended by zero to the common
ambient graph.  Thus a lower bound for a confined family controls a wider
crossing family only when this subwalk property holds.

We specify the spatial tests to which the estimates will apply.  At scale
\(m\) letters per time unit, write \(\Psi_m\) for the reference-coordinate
homeomorphism of Proposition~\ref{prop:geometry}(iii), with \(m=2n\) on a
duration-one sphere.  For a bilateral test, \(\Psi_m\) denotes the
restriction to the unchanged protected flag subcomplex of the homeomorphism
from the finite sewn sphere to the reference one-point compactification
supplied by \cite[Proposition~4.9]{C}.  Edges are drawn in this surface as
in Lemma~\ref{lem:dual-geometry}.  Given open sets
\(P_-,P_+\Subset W\) with disjoint closures, let
\(\mathcal F_m(P_-,P_+;W)\) contain all vertex-to-vertex walks in the chosen
network whose initial and terminal vertices project into \(P_-\) and
\(P_+\), respectively, and whose entire edge arcs project into \(W\).
Every traversed edge is charged in full, including an edge crossing a
boundary of a smaller spatial band.

For \(A(z;r,R)=\{w:r<\abs{w-z}<R\}\), fix
\[
 \mathbf a=(a_0,a_1,a_2,a_3,a_4,a_5),
 \qquad 0<a_0<a_1<a_2<a_3<a_4<a_5.
\]
The annular family \(\mathcal A_m(z,r;\mathbf a)\) is
\(\mathcal F_m(P_-,P_+;W)\) with
\begin{equation}\label{eq:annular-template}
 P_-=A(z;a_1r,a_2r),\qquad
 P_+=A(z;a_3r,a_4r),\qquad
 W=A(z;a_0r,a_5r).
\end{equation}
The positive gaps are part of the template.  Once the flag mesh is smaller
than those gaps, any walk advancing from \(\overline B(z,a_1r)\) to the
exterior of \(B(z,a_4r)\) contains a walk between slightly inset versions
of these ports, confined in \(W\).  Indeed, take the portion after the last
visit to an intermediate inner radius and before the first subsequent visit
to an intermediate outer radius.  The mesh places its endpoint vertices in
the prescribed bands and keeps all its edge arcs in \(W\).

\begin{theorem}[Electrical bounds for fixed tests]\label{thm:network-bounds}
Consider either the primal or the dual network.  Fix a countable collection
of intrinsic instructions from Proposition~\ref{prop:readout-extension}.
Suppose their activated spatial interpretations have the strict ports,
confinements, and complete stars described below, compactly in an ordinary
reference chart.  The same assertions hold in the sphere patches covered
by Proposition~\ref{prop:sphere-support}.  In every joint representation
extraction, almost surely on activation, each fixed instruction has the
following properties for all sufficiently large discrete indices.
\begin{enumerate}[label=\textup{(\roman*)}]
\item For the annular template \eqref{eq:annular-template},
\[
 \EL\bigl(\mathcal A_m(z,r;\mathbf a)\bigr)\ge c_{\mathbf a}>0.
\]
The constant is deterministic and depends only on the relative radii and
the fixed protection margins.
\item Fix \(0<a<b\), with a protected disk buffer around \(B(z,br)\).
There is a function \(u_m\) on the network vertices with values in
\([0,1]\), equal to \(1\) at every vertex whose projected complete flag
star meets \(\overline B(z,ar)\), and equal to \(0\) at every vertex whose
projected star meets
the exterior of \(B(z,br)\).  Its nonzero edge differences are supported
in a fixed compact subannulus of \(A(z;ar,br)\), and
\[
 \mathcal E^0(u_m)\le C.
\]
The constant depends only on \(a,b\) and the chosen relative margins.
In a protected patch the function is extended by the indicated constants
across its two sides.
\item Fix \(0<a<b<c\).  There is a probability law on simple circuits
\(\Gamma_m\) in the network, compactly confined in \(A(z;br,cr)\) and
surrounding \(\overline B(z,ar)\), such that
\[
 \sum_e\PP(e\in\Gamma_m)^2\le C_{a,b,c}.
\]
A loop and a two-edge Jordan circuit are permitted simple circuits.
The constant also records the fixed inner and outer collar margins.
\item Let \(W\) be connected and open, and let \(P_0,P_1\) be nonempty
open subsets of \(W\).  For any fixed compact polygonal route in \(W\)
from \(P_0\) to \(P_1\), there is a probability law \(\mu_m\) of graph
paths between the vertex ports, compactly confined in \(W\), with
\[
 \mathcal T(\mu_m)\le C<\infty .
\]
Here \(C\) may depend on the route, its clearance from \(W^c\), and its
endpoint clearances inside the ports.
\end{enumerate}
For translates and dilates of one fixed relative template, the constants
are the same; they do not depend on the absolute radius.  The eventual
index belongs to the fixed activated test and may depend on it and on the
sample.  The assertions hold simultaneously for the chosen countable
collection.
\end{theorem}

We first prove the lower annular bound.  Lemma~\ref{lem:finite-el} will
then convert it to a cutoff bound, and the common planar drawing will turn
dual cutoffs into primal circuit laws.

\subsection{Annular barriers at power scales}

Fix one of the two networks for this part of the argument.  At diffusive
height scale \(2^j\), use an unconditioned contour coupling with
\(m_j=4^j\) letters per time unit, together with its continuum surface and
protected reference charts.  These are the unconditioned and protected
versions of the encoding and topology in \cite[Sections~3--5]{C}.
Choose a round traversal band in an ordinary part of a cone chart, with its
field buffers strictly inside the normalization circle and away from the
distinguished center.  Fix its center
\(z_0\) and radii
\[
 0<b_0<b_1<\cdots<b_{11}.
\]
Write \(A_0(b,c)=A(z_0;b,c)\), and choose these radii and the field buffers
so that, after one fixed affine change, their closures fit strictly in the
shell and buffers used by \cite[Theorem~6.4]{C}.  Define the middle family
at any effective unit \(m\) by
\[
 \mathcal F_m^{\mathrm{mid}}
 =\mathcal F_m\bigl(A_0(b_1,b_4),A_0(b_7,b_{10});
                         A_0(b_0,b_{11})\bigr),
\]
and let \(Z_j\) be its raw edge extremal length at \(m=m_j\).
The ports have positive thickness and separated closures.  If the
temporary correspondence has not yet reached the required mesh, or if
these ports do not yet give a nonempty family, set \(Z_j=1\).  The
probability of this convention tends to zero.  To see the path-existence
assertion, follow a fixed radial arc with positive clearance from the
confinement boundary.  The faces met along the arc have vanishing diameter;
their boundaries supply a walk in the chosen network, still in the
confinement, once the mesh is small.  The same argument applies in the
other network because complete primal and dual stars have vanishing mesh.
The local finite sewing retains all these incidences.  When the convention
is not used, \eqref{eq:port-resistance} gives \(0<Z_j<\infty\).

The comparison with local numerical observations uses the capture ports
\(A_0(b_2,b_3)\), \(A_0(b_8,b_9)\) and capture confinement
\(A_0(b_1,b_{10})\).  Cover their complete time preimages and stars by
finitely many cells whose outer port and confinement covers still lie in
the corresponding middle sets, as in Section~\ref{sec:locality}.
One activated instruction \(\alpha\) then gives a family
\(\mathcal L_m^\alpha\) satisfying, for all large indices,
\begin{equation}\label{eq:band-sandwich}
 \mathcal L_m^\alpha\subset\mathcal F_m^{\mathrm{mid}},
 \qquad
 \text{every wider traversal has a subwalk in }\mathcal L_m^\alpha .
\end{equation}
Here a wider traversal goes from \(\overline B(z_0,b_1)\) to the exterior
of \(B(z_0,b_{10})\).  Cut it after its last visit to radius
\((b_2+b_3)/2\) and before its first subsequent visit to radius
\((b_8+b_9)/2\).  Vanishing mesh puts its endpoint vertices in the capture
ports and all its arcs in the capture confinement, so the cell covers admit
this subwalk.  This holds for every traversal at once, since the cells
cover full preimages, not chosen paths.
Thus \eqref{eq:el-subpassage} gives
\(\EL(\text{wider traversals})\ge\EL(\mathcal L_{m_j}^\alpha)\ge Z_j\)
at \(m=m_j\), whenever the temporary default is not used.
In the other direction, cut each middle walk after its last visit to radius
\(b_5\) and before its first subsequent visit to radius \(b_6\).
For all large indices the resulting vertex subwalk lies compactly in
\(A_0(b_4,b_7)\) and has endpoint separation at least
\(d_0=(b_6-b_5)/2>0\).  Trimming decreases every edge count.  An outer
cell cover of all these trimmed paths, with port margins around \(b_5\)
and \(b_6\) and confinement still inside \(A_0(b_4,b_7)\), supplies a
local witness for an upper bound on their squared traffic.

The cell-described families \(\mathcal L_m^\alpha\) and the outer covers
of trimmed paths are functions of the canonical finite incidence records
\(\mathcal R_m^\alpha\).  The sets of instructions fitting the
bands with the stated inclusions are chosen from the limiting local input.
They form a countable union of activated descriptions.  In particular,
the extremal length defined using a temporary correspondence is only used
to witness inequalities for these local observations.

Choose a sufficiently small \(\varepsilon>0\), and choose deterministic
lower quantiles \(\rho_j>0\) satisfying
\begin{equation}\label{eq:raw-el-quantiles}
 \PP(Z_j\ge\rho_j/2)\ge1-2\varepsilon,
 \qquad
 \PP(Z_j\le2\rho_j)\ge\varepsilon/2 .
\end{equation}
Such choices exist for positive finite random variables, including laws
with atoms.  We will choose \(\varepsilon\) small enough for the fixed
changes of reference law below.

Suppose for a contradiction that \(\liminf_j\rho_j=0\).  There are running
record indices \(j_k\) such that
\begin{equation}\label{eq:record-low-units}
 m_k=4^{j_k},\qquad \rho_{j_k}\longrightarrow0,\qquad
 \rho_{j_k-i}\ge\rho_{j_k}
 \quad\text{eventually for every fixed }i\ge0.
\end{equation}
Retain the rows with effective units \(m_k4^{-i}=4^{j_k-i}\) for every
fixed backward offset \(i\), and divide all their edge-EL observations by
the same number \(\rho_{j_k}\).  Each fixed effective unit tends to
infinity, as required by Proposition~\ref{prop:readout-extension}.
The first event in \eqref{eq:raw-el-quantiles} in row \(j_k-i\) now gives,
through \eqref{eq:band-sandwich}, a local shell lower bound of the form
\(c\rho_{j_k}\), with a deterministic \(c>0\).

Here is the probability statement about this shell event that will be used.
In each fixed row retain the high-event indicator in
\eqref{eq:raw-el-quantiles}, the continuum input, and the countable local
outputs in units \(\rho_{j_k}\).  Retain also the indicator that
\(Z_{j_k-i}\) used the temporary default.  Its probability tends to zero,
so its Bernoulli limit is zero almost surely and the default is absent
eventually in the representation of each fixed row.
For any instruction \(\alpha\) specifying
a confined port family \(\mathcal F_m^\alpha\), write
\[
 Y_i^\alpha
 =\lim_{k\to\infty}
   \frac{\EL(\mathcal F_{m_k4^{-i}}^\alpha)}{\rho_{j_k}}
 \quad\hbox{in }[0,\infty]
\]
for its jointly extracted output, interpreted through compactification.
Choose a fixed \(c_0>0\) smaller than the lower bound left after all the
band and numerical losses.  Let \(\mathsf A_{\mathrm{sh}}^\alpha\) be the
local activation event that \(\alpha\) has margins in the unit shell
guaranteeing the two relations in \eqref{eq:band-sandwich}.  Define the
unit event
\begin{equation}\label{eq:high-unit-event}
 \mathcal G_i
 =\bigcup_{\alpha}
   \bigl(\mathsf A_{\mathrm{sh}}^\alpha\cap\{Y_i^\alpha>c_0\}\bigr).
\end{equation}
The union is countable and its activation conditions depend only on the
limiting local input.  In a joint representation extraction the limiting
high indicator has probability at least \(1-2\varepsilon\).  On the event
that it equals \(1\), the discrete indicators equal \(1\) eventually and
the temporary default is absent.
A finite cell cover chosen from the limiting full preimages then gives
an activated instruction in \eqref{eq:band-sandwich}; its limiting output
is strictly larger than \(c_0\).  Hence
\(\PP_{\mathrm{row}}(\mathcal G_i)\ge1-2\varepsilon\) for every fixed
\(i\).  This proves a probability statement about the extracted local
array.  It uses neither continuity of \(Z_j\) in the contour nor a
conditional finite-scale probability for activation.

Normalize each row in its own effective unit \(m_k4^{-i}\), using the same
output-independent rule for the cone embedding and stationary imaginary
lift, and take the countable-offset extraction against one fixed
unit-reference input.  The ordinary cone restriction and the unit-shell
restriction of \cite[Theorem~6.4]{C} are locally mutually absolutely continuous,
with the same conditional observation kernel by
Proposition~\ref{prop:readout-extension} and
\cite[Proposition~5.12]{C}.  After the row has been normalized, the input
density \(D\) is fixed, independent of \(i\) and of \(\rho_{j_k}\).
Writing \(\mathsf P_{\mathrm{row}}\) for this common normalized cone input
law, any event \(B\) of the input and its outputs satisfies
\[
 \PP_{\mathrm{ref},i}(B)
 \le L\PP_{\mathrm{row},i}(B)
      +\int_{\{D>L\}}D\,\dd\mathsf P_{\mathrm{row}}
 \qquad(L>0).
\]
This follows by splitting the input density and retaining the common
conditional kernel.  Integrability of \(D\) gives uniform absolute
continuity, independently of the output kernel and denominator.  We may therefore
choose \(\varepsilon\) so that
\(\PP_{\mathrm{ref},i}(\mathcal G_i)\ge1-\varepsilon_*\) for every fixed
\(i\), where \(\varepsilon_*\) is the constant in the aligned-annulus
theorem.  The reference package includes
the imaginary field with its stationary phase and the actual relative
traversal order.  Every label here uses the traversal family for one fixed
compact shell protection, one bounded real buffer around it, and one fixed
larger imaginary buffer.  This family is finite almost surely under one
fixed unclocked imaginary marginal.  The all-order count and weight
exhaustion in Section~6 of~\cite{C} first restricts, with probability loss
\(\zeta\), to at most \(J_0<\infty\) feasible orders whose weights are all
at least \(\delta>0\).  An actual-order failure bound
\(\varepsilon_{\mathrm{ord}}\) then gives all-order failure at most
\(\zeta+\varepsilon_{\mathrm{ord}}/\delta\), uniformly in the labels;
\(\zeta\) is chosen first and then \(\varepsilon_{\mathrm{ord}}\).
For its numerical observations, the proof of
\cite[Theorem~6.4]{C} uses only the local observation kernel, its product
property on disjoint protected supports, and protection for every relative
traversal order.  Proposition
\ref{prop:readout-extension} supplies these properties for the added
edge array, so that proof applies to the local unit events
\(\mathcal G_i\).  We use this extended application below.
The same choice of \(\varepsilon\) may accommodate the
bounded family of constant field shifts used below: their cutoff
Cameron--Martin densities have a uniformly bounded \(L^p\) norm for some
\(p>1\), so Hölder's inequality makes the loss uniformly small.

For a fixed integer \(N\), take
\[
 I_N=[N,2N]\cap\mathbb Z,\qquad v_i=4^{-i}\quad(i\in I_N).
\]
Use \(\mathcal G_{N,i}=\mathcal G_i\) for \(i\in I_N\).  The event with
label \(v_i\) is evaluated in the larger \(m_k\) array by
the indexed field and volume identity \eqref{eq:indexed-units}.  In the
notation of the aligned theorem its normalized real field at center \(c\)
and radius \(r\) is
\begin{equation}\label{eq:aligned-unit-field}
 h^{c,r,v_i}(u)
   =h(c+ru)+Q\log r-\gamma^{-1}\log v_i .
\end{equation}
The raw edges and their EL have not changed, and the denominator is still
\(\rho_{j_k}\).  All offsets in \(I_N\) are fixed before \(k\to\infty\).
Theorem~6.4 of \cite{C} now supplies, almost surely in each ordinary chart
for all sufficiently large \(N\), at most \(CN\) radii in
\([e^{-BN},e^{-aN}]\), with deterministic \(0<a<B\).  At a given radius
\(r\), centers lie on a grid of spacing comparable to \(r\), with bounded
overlap of their buffered shells.  Every point of a prescribed compact
set belongs with strict clearance to the central holes of at least \(cN\)
successful shells at distinct radii.  One successful test is retained per
center and radius.  The constants depend only on the fixed shell buffers
and probability cutoff.  The probability-one statement, with its constants,
persists under fixed locally absolutely continuous changes of the field
law.  A fixed \(N\) gives finitely many tests in a compact set, so their
strict numerical inequalities and incidences hold in the discrete
extraction for all sufficiently large \(k\).

For the next estimates, regard each traversal of an edge as a traversal
of a unit interval.  For a parameter interval \(J\) of a walk, define
\(N_e(\pi;J)\) as the total length of its portions in the unit intervals
corresponding to \(e\), and put
\[
 L_g(\pi;J)=\sum_e g_eN_e(\pi;J).
\]
For a full vertex walk this agrees with \(L_g(\pi)\).  For disjoint
parameter intervals these charges add, even when their images overlap.
In particular, their sum is at most the charge of the whole walk.

\begin{lemma}[Edge density budget]\label{lem:edge-density-budget}
In the record-low extraction \eqref{eq:record-low-units}, let \(U\) be a
finite union of rational boxes compactly in an ordinary reference chart,
and write \(U^{+\delta}\) for its Euclidean \(\delta\)-neighborhood.
Set \(R_N=e^{-aN}\).  Almost surely, for every sufficiently large \(N\)
from the aligned configuration there are nonnegative edge functions
\(g_{k,N,U}\), supported on edges whose entire images lie in
\(U^{+CR_N}\), such that
\begin{align}
 \limsup_{k\to\infty}\rho_{j_k}\norm{g_{k,N,U}}_2^2
     &\le C\,\operatorname{Area}(U^{+CR_N}),\label{eq:edge-budget-energy}\\
 L_{g_{k,N,U}}(\pi;J)
     &\ge c\,\diam \Psi_{m_k}(\pi(J))-CR_N-\epsilon_{k,N,U}.
       \label{eq:edge-budget-charge}
\end{align}
Here \(\epsilon_{k,N,U}\to0\) for fixed \(N,U\), uniformly over every
walk interval whose projected image is in \(U\).  The same density works
for all such intervals and for every finite disjoint collection of them.
The constants \(c,C\) are deterministic and independent of \(k,N,U\)
once the shell geometry is fixed.  The discrete threshold may depend on
the finite configuration and its margins.
\end{lemma}

\begin{proof}
On a successful shell at center \(c\) and radius \(r\), the local
normalized EL lower bound and \eqref{eq:edge-el} give an edge density
\(g_k^{c,r}\), with fixed slack, for which
\begin{equation}\label{eq:one-shell-density}
 L_{g_k^{c,r}}(\pi)\ge r
 \quad\hbox{on every complete shell traversal},\qquad
 \norm{g_k^{c,r}}_2^2\le \frac{Cr^2}{\rho_{j_k}}.
\end{equation}
Its support lies in the shell buffer.  The subwalk inclusion
\eqref{eq:band-sandwich} makes the charge valid for every traversal.
The densities may be chosen measurably: a port family may be restricted
to its finitely many simple paths, and a first rational density satisfying
a strict near-minimum bound gives a choice from the finite incidence
record.

At each radius retain the successful centers whose buffers meet \(U\),
and define
\[
 g_{k,N,U}=\frac1N\sum_{(c,r)}g_k^{c,r}.
\]
For fixed \(N\), the edge mesh is eventually smaller than every shell
margin in this finite collection.  A whole supported edge can then be
assigned to a point in the slightly enlarged buffer that contains it.
The bounded spatial overlap at one radius bounds the number of summands
on this edge by a deterministic constant; no degree bound is used.
Grid counting and \eqref{eq:one-shell-density} give
\[
 \norm{\sum_c g_k^{c,r}}_2^2
 \le C\sum_c\norm{g_k^{c,r}}_2^2
 \le \frac{C\,\operatorname{Area}(U^{+CR_N})}{\rho_{j_k}}.
\]
There are at most \(CN\) radii.  Cauchy--Schwarz for their sum, followed
by division by \(N^2\), proves \eqref{eq:edge-budget-energy} and the
support assertion.

For the charge, choose a coordinate whose oscillation on \(\Psi_{m_k}(\pi(J))\)
is at least its diameter divided by \(\sqrt2\), and use a subinterval
between a minimum and a maximum in that coordinate.  Ignore coordinate
levels within \(CR_N\) of the two ends.  At each remaining level, a visit
of the walk lies in at least \(cN\) successful central holes.  A hole of
radius \(r\) accounts for a coordinate interval of length at most \(Cr\).
For every hole counted in this way the path makes a complete traversal of
the corresponding buffered shell.  More precisely, strict clearance and
the fixed-\(N\) mesh bound place a full vertex-to-vertex shell subwalk
inside \(J\), even if an endpoint of \(J\) is in the interior of an edge.
Thus each counted shell contributes its full charge in
\eqref{eq:one-shell-density}.  Integrating over the coordinate levels gives
\[
 cN\bigl(\operatorname{osc}-CR_N-\epsilon_{k,N,U}\bigr)
 \le C\sum_{\text{distinct shells traversed}}r .
\]
The summed density charges at least the sum on the right before the fixed
constant.  Divide by \(N\) and absorb constants.  This proves
\eqref{eq:edge-budget-charge}.  Its error is uniform because all projection
and edge-mesh errors are uniform for this fixed configuration.  Charges
of disjoint parameter intervals add by their fractional edge definition;
no endpoint edge is rounded up separately.  This is the level-integration
argument of \cite[Section~8]{C} applied to unit edge intervals.
\end{proof}

\subsection{Rectifiable limits of inexpensive crossings}

We record how the budget controls path laws without a bound on the
Euclidean length of an individual projected edge.  The estimate is for
the average arclength of the limiting paths.

\begin{lemma}[Compactness of edge-path laws]\label{lem:edge-flow-compactness}
Fix a sample on which the budgets of Lemma~\ref{lem:edge-density-budget}
hold for all sufficiently large members of a deterministic sequence of
\(N\)'s and the countable box library.  Suppose \(\mu_k\) are laws of
paths confined in one compact set \(K\) in the chart, their projected
endpoints are separated by at least \(d_0>0\), and
\begin{equation}\label{eq:cheap-edge-flow}
 \mathcal T(\mu_k)\le C_0\rho_{j_k}.
\end{equation}
After reparameterization and a subsequence, their projected laws converge
weakly on \(C([0,1],K)\) to a probability law \(\mu\).  Its paths retain
endpoint separation \(d_0\) and are rectifiable almost surely.  Their
average arclength measure \(\Lambda\), counted with multiplicity, satisfies
\begin{equation}\label{eq:average-edge-arclength}
 \Lambda(U)\le C\sqrt{\operatorname{Area}(U)}
\end{equation}
for every finite union of rational boxes compactly in the chart.
Consequently \(\Lambda\) is nonzero, finite on \(K\), and absolutely
continuous with respect to planar area.
\end{lemma}

\begin{proof}
For every edge density, the traffic pairing is
\begin{equation}\label{eq:traffic-pairing}
 \EE_{\mu_k}L_g(\pi)
   =\sum_e g_e\tau_{\mu_k}(e)
   \le\norm{g}_2\sqrt{\mathcal T(\mu_k)}.
\end{equation}
For \(g=g_{k,N,U}\), equations \eqref{eq:edge-budget-energy} and
\eqref{eq:cheap-edge-flow} bound this expectation by
\(C\sqrt{\operatorname{Area}(U^{+CR_N})}\) for large \(k\).

Take \(U\) large enough to contain \(K\) with room.  Given \(d>0\), choose
a fixed successful \(N\) with \(CR_N<cd/16\), and then take \(k\) large
enough that \(\epsilon_{k,N,U}<cd/16\).  Every path interval of projected
diameter greater than \(d/2\) pays at least \(cd/4\) in
\eqref{eq:edge-budget-charge}.  By continuity, the endpoints of any finite
collection of intervals of diameter at least \(d\), with pairwise disjoint
interiors, can be approximated from within by rational endpoints while
retaining diameter greater than \(d/2\).  The supremum of the crossing
counts over such rational collections is measurable, since the collections
are countable.  Charges add, so the simultaneous bound and
\eqref{eq:traffic-pairing} bound its expectation by \(C/d\); formally,
increase over finitely many rational collections and use monotone
convergence.  This also bounds the maximum count for the original
intervals.

For completeness, a measurable clock gives the required tightness.
First normalize the original parameter of a continuous projected path to
\([0,1]\).  Mark its successive first moves by
distance \(2^{-\ell}\) and let \(J_\ell\) be their number.  Give each such
mark mass \(2^{-\ell}/(1+J_\ell)\), for every \(\ell\ge1\), and add
Lebesgue mass in the original parameter.  The total mass is between \(1\)
and \(2\).  Its inverse clock, with a pause at each atom and normalization
to \([0,1]\), represents the same path.  On a set where
\(J_\ell\le M_\ell\) for every \(\ell\), a parameter interval shorter than
\(2^{-\ell-1}/(1+M_\ell)\) cannot pass through a complete intervening
pause; its image oscillation is at most \(2^{2-\ell}\).
The crossing-count bounds permit \(M_\ell<\infty\) to be chosen so that
the total exceptional probability is arbitrarily small.  For each
\(\ell\), the finitely many indices preceding its eventual bound are
included by increasing \(M_\ell\).  Arzel\`a--Ascoli and Prokhorov's
theorem give tightness and a weak limit on \(C([0,1],K)\).  The endpoint
separation passes to this limit.

To control variation, fix an integer \(J\) and a finite rational box
union \(U\).  Let \(V_{J,U}(\gamma)\) be the supremum of the sums
\[
 \sum_{\ell=1}^j\abs{\gamma(t_\ell)-\gamma(s_\ell)},\qquad 0\le j\le J,
\]
over disjoint closed intervals with rational endpoints whose images have
positive clearance inside \(U\), with the empty sum equal to \(0\).
This is lower semicontinuous on path
space: each strict containment is open and each corresponding finite sum
is continuous there.  The simultaneous charge bound, applied before
taking this supremum, and \eqref{eq:traffic-pairing} yield
\[
 \EE_{\mu_k}V_{J,U}
 \le C\sqrt{\operatorname{Area}(U^{+CR_N})}
       +CJR_N+J\epsilon_{k,N,U}.
\]
First pass to the weak path-law limit \(k\to\infty\), with \(N,J\) fixed.
Then let \(N\to\infty\) through successful configurations, keeping \(J\)
fixed, and finally let \(J\to\infty\).  The finite interval sums exhaust
variation in \(U\).  Applying this first to a larger confinement proves
finite total variation for almost every path, and applying it in each
\(U\) gives \eqref{eq:average-edge-arclength}.  Each path has length at
least \(d_0\), so \(\Lambda\ne0\).  Covering a planar null set by rational
boxes of arbitrarily small total area, and increasing over finite
subunions, proves \(\Lambda\ll\dd z\).  This proves the edge version of
the flow argument in \cite[Section~8]{C}.
\end{proof}

The preceding lemma may be applied sample by sample on a positive-probability
event of a joint extraction.  One first restricts to finite bounds for the
confinement and traffic and retains the path laws along the tight
subsequence.  These laws serve to witness events in the local EL arrays;
the conditional kernels of Section~\ref{sec:locality} are not conditioned
on a path law.

\subsection{A primal--dual obstruction to degeneration}

\begin{proposition}[Positive raw crossing quantiles]\label{prop:edge-quantiles}
For either choice of network in the power-scale construction,
\[
 \liminf_{j\to\infty}\rho_j>0.
\]
The lower bound is deterministic.
\end{proposition}

\begin{proof}
Continue the record-low assumption \eqref{eq:record-low-units}.  Choose the
joint extraction used above to retain, in addition, the duration-one
spheres with \(n_k=m_k/2\), their rooted dual experiments, the matching
ordinary reference packages for the two sphere marginals, and the cone
packages witnessing the quantile events for each reference marginal.
Propositions~\ref{prop:readout-extension} and \ref{prop:sphere-support}
provide these matched packages, and Lemma~\ref{lem:dual-geometry} couples
the two sphere marginals on the same physical maps with the same edge
count.  The high-unit and budget arguments above apply to each reference
marginal.  We carry out the germ argument below in each one, using its own
conditional kernel for both its cone transfer and its later sphere
transfer; the notation first describes one such matched pair.

Retain also the low-event indicators in \eqref{eq:raw-el-quantiles}.  Their
limiting event has positive probability.  On that event the discrete
indicators equal \(1\) eventually.  Since \(\rho_{j_k}\to0\), the default
value \(Z_{j_k}=1\) cannot satisfy the low inequality for large \(k\).
Lemma~\ref{lem:finite-el} then gives
near-minimum path laws for the middle family, and trimming between the
narrower bands gives squared traffic at most \(C\rho_{j_k}\), endpoint
separation at least \(d_0\), and compact confinement.  A finite outer
cell cover can admit every trimmed path while keeping these two geometric
properties, with strict margins.  Its finite path family contains the
trimmed laws for all large \(k\), so its EL divided by \(\rho_{j_k}\)
is bounded by a fixed finite constant.  Let
\(\mathsf A_{\mathrm{out}}^\alpha\) be the local activation event that
\(\alpha\) is such an outer cover with the stated strict margins.  Thus
the low indicator witnesses the local event
\begin{equation}\label{eq:low-unit-event}
 \mathcal G_{\mathrm{low}}
 =\bigcup_{\alpha}\ \bigcup_{M\ge1}
  \bigl(\mathsf A_{\mathrm{out}}^\alpha\cap\{Y_0^\alpha<M\}\bigr).
\end{equation}
The union is countable and has positive probability in the cone
restriction.  By countability and compact exhaustion of activation,
one instruction \(\alpha\), one finite \(M\), and fixed compact
confinement and positive endpoint separation can be selected on a
positive-probability subevent.  This is an event of the local input and
its normalized numerical output.  Local equivalence of the cone and
ordinary field restrictions, with the common observation kernel,
transfers this event with positive probability to its matching ordinary
reference marginal.
There the convergence \(Y_0^\alpha<M\) and exact reconstruction on
activation give finite normalized EL at most \(M+1\) for all sufficiently
large \(k\).  Lemma~\ref{lem:finite-el} supplies discrete path laws in
this one local family with squared traffic at most
\((M+2)\rho_{j_k}\).  These laws now live on the same ordinary extraction
as the probability-one budgets of Lemma~\ref{lem:edge-density-budget}.
Apply Lemma~\ref{lem:edge-flow-compactness} on their common event to obtain
a law \(\mu\) of nonconstant rectifiable paths and a nonzero measure
\(\Lambda\ll\dd z\).

We describe the local event witnessed by their tangents.  Let \(v\) be a
unit direction and \(v^\perp\) its quarter-turn.  For \(z\in\C\) and
\(s>0\), set
\[
 Q_v(z,s)=
 \{z+tv+uv^\perp:\abs t<3s/2,\ \abs u<s/16\}.
\]
A strict crossing has a compact subroute in \(Q_v(z,s)\), from the portion
with \(t<-s\) to the portion with \(t>s\), with its two vertex ports and
all buffered confinement described by intermediate cells.  For an
instruction \(\alpha\) of this kind, write \(Y^\alpha=Y_0^\alpha\) for
the limiting EL of its path family divided by \(\rho_{j_k}\).  Define
\begin{equation}\label{eq:crossing-germ}
 E_v(z)=
 \bigcup_{\ell_0\ge1}\ \bigcap_{\ell\ge\ell_0}\
 \bigcup_{\alpha}\ \bigcup_{M\ge1}
 \left\{
 \begin{array}{l}
 \alpha\text{ is activated as a strict crossing}\\[-2pt]
 \text{with buffered support in }Q_v(z,2^{-\ell}),\
 Y^\alpha<M
 \end{array}\right\}.
\end{equation}
The unions use the countable instruction list and integer \(M\).  Whether
an instruction fits the displayed open tube is a measurable condition
on its limiting local input.  Thus \(E_v(z)\) is jointly measurable.
Deleting finitely many radii leaves it unchanged, and all supports shrink
to \(z\).  Its indicator is therefore a germ in
Lemma~\ref{lem:deterministic-germs}.  The finite bound \(M\) may depend on
the radius.

At arclength-almost every interior visit of almost every path of \(\mu\),
arclength parameterization is differentiable with a unit tangent.  Choose
a direction from a fixed finite sufficiently fine net so that it lies
within the angular slack of this tangent.  Differentiability gives a
strict crossing of \(Q_v(z,s)\) at every sufficiently small \(s\).
Almost every path belongs to the support of \(\mu\).  For each one of
these fixed strict crossings, choose a smaller compact subroute with
additional margin.  Choose nested open vertex ports
\(P_\pm^{\mathrm{in}}\Subset P_\pm^{\mathrm{out}}\) around its two ends
and nested confinements
\(W^{\mathrm{in}}\Subset W^{\mathrm{out}}\Subset Q_v(z,s)\) around the
subroute; the outer ports still lie beyond the two longitudinal levels.
Cover the full preimages of the inner sets by cells lying in the outer
sets.  This gives one activated outer instruction.  An open neighborhood
\(\mathcal O\) of the path can require the corresponding subroute to
visit the inner ports and remain in the inner confinement.  Full-preimage
reconstruction and the vanishing mesh then imply that, for all large
\(k\), every discrete path with projection in \(\mathcal O\) has a trimmed
vertex subroute admitted by this fixed instruction.  Since the path is
in the support, \(\mu(\mathcal O)>0\).  Weak convergence gives mass at
least \(p>0\) in \(\mathcal O\) for all sufficiently large \(k\).
Restricting and normalizing that mass, and trimming its paths, bounds
traffic coordinatewise by
\(\tau_{\mu_k}/p\).  Its squared traffic is at most
\(C_0\rho_{j_k}/p^2\).  Lemma~\ref{lem:finite-el} proves \(Y^\alpha<\infty\)
for this radius.  This proves \(E_v(z)\) at \(\Lambda\)-almost every
such visit for some direction in the finite net.

Since \(\Lambda\) is nonzero and absolutely continuous, on the
positive-probability flow event at least one of these finitely many
directions satisfies \(E_v\) on a set of positive planar area.
Take here the library of all finite intrinsic instructions and its common
refinements.  Translation by \(a\) carries the fitting condition for
\(E_v(z)\) to that for \(E_v(z+a)\); rotation by \(R\) carries it to that
for \(E_{Rv}(Rz)\).  The same intrinsic instructions describe the
transformed cells, and their edge counts and normalized outputs are
unchanged.  The buffered supports transform with them.  These actions are
therefore exact on the local arrays, while the corresponding ordinary
field laws are locally equivalent as in
Lemma~\ref{lem:deterministic-germs}.  That lemma
therefore makes the indicator for each direction one deterministic
constant almost everywhere on an ordinary connected chart.  The
positive-area event forces the constant for one direction to be \(1\).
Rotational covariance then gives constant \(1\) for every member of a
chosen countable dense direction list.  Countable intersection and an
ordinary chart exhaustion show that, almost surely,
\begin{equation}\label{eq:ubiquitous-edge-germs}
 E_v(z)\text{ holds for planar-area-almost every }z
 \text{ and every direction }v\text{ in this list}.
\end{equation}
This assertion has so far been proved for the one network fixed at the
start of the power-scale construction.

We use the sphere components retained in this joint extraction to obtain
a transverse obstruction.  For a rational compact disk \(K\) in a buffered ordinary
patch and a direction in the countable list, failure is the measurable event
\[
 \left\{\int_K\bigl(1-\ind_{E_v(z)}\bigr)\dd z>0\right\}.
\]
Deleting a common finite number of radii for \(z\in K\) computes the tail
in \eqref{eq:crossing-germ} using only supports inside this patch.  The
failure event is therefore measurable from the local input and the
activated array of normalized EL observations in
Proposition~\ref{prop:sphere-support}, and it is null by
\eqref{eq:ubiquitous-edge-germs} in each matching reference marginal,
with that marginal's own conditional kernel.
Proposition~\ref{prop:sphere-support} transfers the event to the
corresponding sphere marginal, with the same denominator
\(\rho_{j_k}\), actual height, and traversal order.  Auxiliary mark pairs
exhaust the patches off the exploration root.

The two retained sphere experiments have the same law and the same number
of letters per unit.  The chosen network in their two marginals gives the
two opposite physical networks.  The transferred conclusions hold for
both simultaneously.  Their flag surfaces are identified, and the
same three fresh flag-area marks give one canonical coordinate after
composing one coordinate with conjugation if its orientation convention
requires this.  Let \(H_k^+\) and \(H_k^-\) be the two reference
homeomorphisms.  The canonical comparison in Proposition~\ref{prop:geometry},
from \cite[Theorem~8.8]{C}, permits a further extraction on which their changes
to the common canonical coordinate converge uniformly to maps \(F^+\)
and \(F^-\), each conformal or anticonformal.

Choose a point \(w\) in a canonical coordinate plane away from the finitely
many omitted marks and poles, at which the two transported versions of
\eqref{eq:ubiquitous-edge-germs} hold.  Such points have full planar area:
the limiting conformal and anticonformal maps are smooth local
diffeomorphisms and preserve area-null sets.  Let \(z^\pm=(F^\pm)^{-1}(w)\).
Choose directions from the two dense lists whose images under
\(DF^\pm(z^\pm)\) are respectively horizontal and vertical to a
sufficiently small angular error.  Write \(b_\pm>0\) for the scale factors
of these similarity or anti-similarity derivatives.  The two germs allow
all sufficiently small dyadic radii.  We may choose such radii \(s_\pm\)
with
\[
 1\le
 \frac{\max\{b_+s_+,b_-s_-\}}{\min\{b_+s_+,b_-s_-\}}
 \le2.
\]
They may also be chosen small enough that, for a prescribed small
\(\varepsilon_{\mathrm{lin}}>0\), each limiting coordinate change satisfies
\[
 \sup_{\abs{\zeta-z^\pm}\le2s_\pm}
 \abs{F^\pm(\zeta)-F^\pm(z^\pm)-DF^\pm(z^\pm)(\zeta-z^\pm)}
 \le \varepsilon_{\mathrm{lin}} b_\pm s_\pm .
\]

Put \(s_*=\min\{b_+s_+,b_-s_-\}\).  If the angular error is \(\delta\),
the transverse deviation of either transformed tube is at most
\[
 2s_*\bigl(1/16+(3/2)\sin\delta+\varepsilon_{\mathrm{lin}}\bigr),
\]
and its endpoint advance in its longitudinal direction is at least
\(s_*(\cos\delta-\sin\delta/16-\varepsilon_{\mathrm{lin}})\).
Choose \(\delta,\varepsilon_{\mathrm{lin}}\) so that
the first bound is less than \(s_*/4\) and the second exceeds \(0.8s_*\).
Thus, in coordinates centered at \(w\), every path in
the first transformed family stays in \(\{\abs y<s_*/4\}\) and reaches
both \(x<-0.8s_*\) and \(x>0.8s_*\); every path in the second stays in
\(\{\abs x<s_*/4\}\) and reaches both \(y<-0.8s_*\) and \(y>0.8s_*\).
Any two such continuous paths intersect, by separation in the intervening
rectangle.  Fix these radii, activated covers, and their finite bounds
from \eqref{eq:crossing-germ}.  Uniform convergence of the coordinate
changes and the strict tube margins give the same forced intersection
for every pair of discrete paths for all sufficiently large \(k\).

The two discrete EL values are at most \(C_+\rho_{j_k}\) and
\(C_-\rho_{j_k}\), with finite constants because the tests have now been
fixed.  Choose near-minimum path laws by Lemma~\ref{lem:finite-el}, and
write their traffics on the paired physical edges as \(\tau_k(e)\) and
\(\tau_k^*(e^*)\).  Every pair of sampled paths intersects.  By
Lemma~\ref{lem:dual-geometry}, a primal and a dual edge meet only when
they are paired, so for every pair of walks
\(\sum_eN_e(\pi)N_{e^*}(\pi^*)\ge1\).  Fix the full coupled record, the
selected families, and both path laws, including their possible dependence
on the continuum data, radii, and orientation sign.  Take the product of
these two fixed laws.  Taking expectations under this pathwise product and applying
Cauchy--Schwarz gives
\[
 1\le\sum_e\tau_k(e)\tau_k^*(e^*)
 \le
 \left(\sum_e\tau_k(e)^2\right)^{1/2}
 \left(\sum_e\tau_k^*(e^*)^2\right)^{1/2}
 \le\sqrt{C_+C_-}\,\rho_{j_k}\longrightarrow0.
\]
This is impossible.  It excludes the record-low extraction for the
chosen network, and the same argument applies when the initial choice
is exchanged.  Since each quantile sequence is deterministic, each
positive lower limit is deterministic.
\end{proof}

\subsection{Bounds at arbitrary units and their consequences}

\begin{proof}[Proof of Theorem~\ref{thm:network-bounds}]
Let \(m_k\to\infty\) be any sequence of letters per time unit.  Write
\[
 m_k=4^{j_k}v_k,\qquad 1\le v_k<4,
\]
and pass further so that \(v_k\to v\in[1,4]\), jointly with all fixed
backward rows.  Proposition~\ref{prop:edge-quantiles} makes the high
barriers uniformly positive in raw EL units in those power rows.
Equation~\eqref{eq:indexed-units} expresses the same incidence record in
the \(m_k\) units: the area clock is divided by \(v_k\), frontier heights
by \(\sqrt{v_k}\), and the real field is shifted by
\(-\gamma^{-1}\log v_k\).  A backward row \(4^{j_k-i}\) has relative label
\(v_k^{-1}4^{-i}\).  Thus the aligned theorem has multiplier \(M=v^{-1}\)
in the limit.  The normalized field identity is, explicitly,
\[
 \bigl(h-\gamma^{-1}\log v\bigr)(c+ru)+Q\log r
       -\gamma^{-1}\log(v^{-1}4^{-i})
 =h(c+ru)+Q\log r+i\gamma^{-1}\log4.
\]
The raw EL is unchanged.  Strict margins allow the converging ratios and
the rounded units to be retained in the indexed arrays.  The bounded
constant shifts have uniformly bounded cutoff Cameron--Martin norms on
the field buffers, as used when choosing \(\varepsilon\) above.  Therefore
the same high unit-reference probability applies before annular alignment.
This is the field and unit change of \cite[Proposition~5.12]{C}; it
re-expresses one raw array and does not compare differently indexed
kernels at an unchanged field.

Apply the aligned theorem again and repeat
Lemma~\ref{lem:edge-density-budget}, now with a fixed positive raw barrier,
so that the denominator in \eqref{eq:edge-budget-energy} is \(1\).
For a fixed annular template, choose slightly inset ports between
\(a_2r\) and \(a_3r\), and an outer cell confinement inside \(W\).
It catches a subwalk of every member of
\(\mathcal A_m(z,r;\mathbf a)\).  Cover that confinement by a finite
rational box union \(U\) of area at most \(C_{\mathbf a}r^2\).  Every
caught subwalk has projected diameter at least \(c_{\mathbf a}r\).
On the probability-one aligned event choose a successful, fixed \(N\)
large enough that \(CR_N<c_{\mathbf a}r/4\) and
\(\operatorname{Area}(U^{+CR_N})\le C_{\mathbf a}r^2\).  For this fixed
configuration, all sufficiently large discrete indices have a density
charging every such subwalk by at least \(c'_{\mathbf a}r\), with squared
norm at most \(C'_{\mathbf a}r^2\).  Equation~\eqref{eq:edge-el} and the
subwalk property give
\[
 \EL\bigl(\mathcal A_m(z,r;\mathbf a)\bigr)
 \ge\frac{(c'_{\mathbf a}r)^2}{C'_{\mathbf a}r^2}
 =c''_{\mathbf a}>0.
\]
The constants are deterministic.  The choice of a successful \(N\) may
depend on the activated sample and the fixed radius, but it is fixed before
the discrete limit.  Narrow collars are included by choosing their
positive relative gaps first.

The passage to \(v_k\to v\) preserves the quantifier over representation
extractions.  Start with any such extraction and retain its existing
compactified local annular output limits, together with their activation
inputs, in the further extraction above.  Their marginal law is preserved.
If an activated limit failed a fixed slack version of the deterministic
lower bound with positive probability, the same failure would have positive
probability in the further extraction, contradicting the bound just proved.
Thus every original limiting output has that lower bound on activation.
Choosing a strictly smaller deterministic constant and using compactified
convergence gives the eventual finite inequality in the original
representation.  A countable intersection treats the chosen library;
the index at which an inequality starts can still depend on its sample
and its fixed test.

The local lower tests in this argument are finite incidence observations
with the band inclusions \eqref{eq:band-sandwich}.  Proposition
\ref{prop:sphere-support} transfers their probability-one limiting bounds
to sphere patches, and convergence of the compactified observations gives
the eventual finite inequalities with a fixed slack.  A finite family of
ordinary patches also treats a collar surrounding an isolated omitted
point: choose inner patches with a Lebesgue number comparable to the
radius and larger patches still avoiding the point.  A through walk has
a subinterval of diameter comparable to the radius in one of the larger
patches.  Summing the finitely many corresponding densities charges all
through walks with the same kind of shape bound.  The number and relative
geometry of these patches can be fixed from the collar ratios.  The
protected enclosure of Proposition~\ref{prop:geometry} keeps the resulting
separation in the original surface.  This proves part (i), including its
stated sphere uses.

We turn this estimate into part (ii).  Fix relative radii
\[
 a<\beta_0<\beta_1<\cdots<\beta_5<b.
\]
In the protected finite graph containing the complete stars over
\(\overline B(z,br)\) and its disk buffer, set
\[
 A_m=\{v:\Psi_m(v)\in\overline B(z,\beta_1r)\},
 \qquad
 B_m=\{v:\Psi_m(v)\notin B(z,\beta_4r)\}.
\]
Part (i) for the template \((\beta_0,\ldots,\beta_5)\) supplies a density
\(g\), normalized to charge its annular family by at least \(1\), with
\(\norm g_2^2\le C\).  Extend it by zero outside its confined edges.
Its support lies compactly in \(A(z;ar,br)\).  Once the mesh is small,
every walk from \(A_m\) to \(B_m\) contains a subwalk of this annular
family and therefore has \(g\)-length at least \(1\).  On the buffered
graph define
\[
 u_m(v)=1-\min\{1,d_g(A_m,v)\},
\]
where \(d_g=\infty\) on components not meeting \(A_m\).  The function is
\(1\) on \(A_m\) and \(0\) on \(B_m\).  As in Lemma~\ref{lem:finite-el},
\(\abs{d_eu_m}\le g_e\), so \(\mathcal E^0(u_m)\le\norm g_2^2\).
Moreover \(u_m\) is constant across every edge on which \(g_e=0\).
The vanishing complete-star mesh and the gaps \(a<\beta_1\),
\(\beta_4<b\) put every vertex whose projected star meets
\(\overline B(z,ar)\) in \(A_m\), and every vertex whose projected star
meets the exterior of \(B(z,br)\) in \(B_m\).  The support of the
differences remains inside the stated compact subannulus.
For a protected use, extend the constants across the interior and
exterior.  Full stars are retained throughout the collar, and any path
from one side to the other contains the tested passage; hence no omitted
incident edge introduces a new difference in this extension.

For part (iii), take the cutoff from part (ii) in the dual network, with
all its nonzero differences in a compact subcollar of \(A(z;br,cr)\).
Its inner constant region contains every face meeting
\(\overline B(z,ar)\), and its outer constant region contains every face
on the other side of the collar.  Choose the subcollar with an additional
fixed margin to both boundary circles.  A paired primal and dual edge
lie in their common flags, whose projected diameters tend to zero by
Proposition~\ref{prop:geometry}; the complete primal and dual star mesh
has the same property.  Thus each primal edge paired with a nonzero dual
difference lies, for all large indices, in a slightly larger compact
subcollar of \(A(z;br,cr)\).  Choose \(t\) uniformly from \((0,1)\).
For all \(t\) except the finitely many vertex values, retain each primal
edge whose two incident dual values are on opposite sides of \(t\), and
orient it along the boundary of the face on the \(u_m>t\) side.  Cyclic
incidence at a primal vertex balances incoming and outgoing edges, so
these edges form closed walks.  A transverse arc from the inner region
to the outer one has total signed transition count of absolute value \(1\).
The closed
walks therefore have nonzero total separating winding.  Splitting at
repeated vertices gives simple circuits, at least one of which surrounds
the inner disk.  All its edges remain in the subcollar.  This reasoning
also includes a primal loop and two parallel edges forming a Jordan
curve; a primal bridge has a dual loop and cannot be a transition edge.

Choose one surrounding circuit by any fixed rule on the finite record.
An edge can be chosen only when \(t\) lies between its two dual values.
Consequently
\[
 \PP(e\in\Gamma_m)\le\abs{d_{e^*}u_m},
 \qquad
 \sum_e\PP(e\in\Gamma_m)^2
 \le\sum_{e^*}(d_{e^*}u_m)^2\le C.
\]
The full face stars and the two constant margins used here are part of
the protection; they ensure a closed separating boundary rather than a
boundary cut off by the edge of a chart.  Exchanging the two networks
proves the circuit assertion in either one.

Finally, take the route in part (iv).  Its compactness gives positive
clearance from \(W^c\), and its endpoints have positive clearance inside
the two ports.  Choose a sufficiently small common radius \(r_0\).
Starting at the first endpoint, choose successive centers at the first
points of the route at Euclidean distance \(r_0\) from the preceding
center.  This terminates after finitely many steps because the polygonal
route has finite length.  When the remaining route stays within that
distance, the last circle can already be chosen wholly inside the terminal
port by the endpoint clearance.  The first circle is wholly in the
initial port, and all circles and their narrow collars lie compactly
inside \(W\).

Use part (iii) for a circuit surrounding \(B(c,3r_0/4)\) and lying in
\(A(c;7r_0/8,9r_0/8)\) at each center \(c\).  Consecutive centers have
distance \(r_0\).  The two enclosed disks of radius \(3r_0/4\) overlap,
whereas each has a point outside the other's disk of radius \(9r_0/8\).
Two disjoint Jordan curves would have
disjoint or nested interiors, so every pair of consecutive circuits
intersects.  In one embedded network an intersection gives a common
vertex or edge.  The union of the sampled circuits is therefore connected
from the first port to the last.  Select a simple graph path in this
union.  If there are \(J_*\) circuits and
\(\tau_i(e)=\PP(e\in\Gamma_i)\), the path traffic is at most
\(\sum_{i=1}^{J_*}\tau_i(e)\) on every edge.
Hence
\[
 \mathcal T(\mu_m)
 \le\norm{\sum_{i=1}^{J_*}\tau_i}_2^2
 \le J_*\sum_{i=1}^{J_*}\norm{\tau_i}_2^2
 \le J_*^2 C.
\]
The number \(J_*\) and the collar shape depend only on the fixed route and
its clearances and are preserved by a common dilation.  This proves
part (iv) and completes the theorem.
\end{proof}
